% Chapter 2 worked example. Included by ch02_resources.tex.
\section{Worked Example: Readmission Analytics Readiness and Data-Foundation Investment}
\label{sec:comp-ch02-worked-example}

\subsection{Purpose and Status}
This worked example demonstrates how the methods in Chapter~2 can be
translated into a reviewable institutional decision. All numerical inputs are
synthetic unless a source is identified explicitly. They are intended to show the
logic of the analysis, not to report empirical findings or establish a universal
threshold. Local use requires replacement of the assumptions with current,
versioned evidence and approval through the relevant clinical, managerial,
economic, legal, technical, and governance processes.

A multi-site hospital group proposes a readmission-risk service, but discharge data, medication reconciliation, laboratory results, and prior admissions are distributed across several systems. The decision is whether to procure a model immediately or first strengthen reusable digital foundations. 

\subsection{Decision Problem and Comparators}
\textbf{Decision problem.} Which sequence of investment offers the greatest institutional value: immediate procurement, a foundation-first programme, or a limited proof of concept restricted to one service?

The comparator set is deliberately broader than the existing pathway. This prevents
a new technology from receiving a lower evidentiary threshold than staffing,
workflow, shared-service, conventional digital, or no-change alternatives.

\begin{longtable}{@{}p{0.08\textwidth}p{0.84\textwidth}@{}}
\caption{Comparators for the Chapter 2 worked example}
\label{tab:comp-ch02-comparators}\\
\toprule
\textbf{Option} & \textbf{Comparator or strategy} \\
\midrule
\endfirsthead
\toprule
\textbf{Option} & \textbf{Comparator or strategy} \\
\midrule
\endhead
\bottomrule
\endlastfoot
1 & Immediate vendor procurement \\
2 & Foundation-first investment in identity, terminology, interfaces, provenance, and monitoring \\
3 & Single-site proof of concept with manual data preparation \\
4 & Conventional discharge-risk checklist and additional care-coordination capacity \\
\end{longtable}

\subsection{Model and Decision Structure}
The analytical structure should follow the decision rather than the available
software. The team should map the causal pathway from intervention input to
information, human decision, action, outcome, resource consequence, and review.
The structure should also identify capacity constraints, uncertainty, subgroup
variation, implementation requirements, and events that would change the
recommendation.

A compact quantitative relationship used in this example is:
\begin{equation}
G=\prod_{k=1}^{K}g_k,\qquad g_k\in\{0,1\}
\label{eq:comp-ch02-key-relationship}
\end{equation}
The equation is an organizing aid rather than a substitute for disaggregated evidence,
qualitative findings, or mandatory safeguards. Its variables and interpretation must
be defined in the local protocol.

\subsection{Synthetic Parameters and Evidence Sources}
Table~\ref{tab:comp-ch02-parameters} provides the base-case inputs. A
production analysis should add parameter distributions, correlations, structural
alternatives, data dates, system versions, and evidence-confidence ratings.

\begin{longtable}{@{}p{0.32\textwidth}p{0.22\textwidth}p{0.38\textwidth}@{}}
\caption{Synthetic parameters for the Chapter 2 worked example}
\label{tab:comp-ch02-parameters}\\
\toprule
\textbf{Parameter} & \textbf{Base value} & \textbf{Source or interpretation} \\
\midrule
\endfirsthead
\toprule
\textbf{Parameter} & \textbf{Base value} & \textbf{Source or interpretation} \\
\midrule
\endhead
\bottomrule
\endlastfoot
Hospitals in scope & 4 & Portfolio plan \\
Interfaces requiring remediation & 11 & Architecture inventory \\
Records with reliable medication reconciliation & 68\% & Data-quality profile \\
Cross-site patient-identity match rate & 96.1\% & Master-patient-index audit \\
Estimated foundation investment & EUR 740,000 & Programme estimate \\
Use cases expected to reuse the foundations & 7 & Digital portfolio \\
Estimated procurement delay from foundation programme & 9 months & Implementation plan \\
\end{longtable}

\subsection{Step-by-Step Analysis}
The sequence below is intended to make the analysis reproducible and to ensure
that evidence, economics, governance, and implementation develop together.

\begin{longtable}{@{}p{0.07\textwidth}p{0.55\textwidth}p{0.30\textwidth}@{}}
\caption{Analytical sequence}
\label{tab:comp-ch02-analysis-steps}\\
\toprule
\textbf{Step} & \textbf{Required work} & \textbf{Decision record} \\
\midrule
\endfirsthead
\toprule
\textbf{Step} & \textbf{Required work} & \textbf{Decision record} \\
\midrule
\endhead
\bottomrule
\endlastfoot
1 & Specify the data, timing, and workflow dependencies of the intended use. & Evidence and responsible owner to be recorded \\
2 & Assess readiness by use case rather than by one maturity score. & Evidence and responsible owner to be recorded \\
3 & Identify which weaknesses affect several proposed services and estimate their portfolio-level reuse value. & Evidence and responsible owner to be recorded \\
4 & Compare immediate procurement with a foundation-first sequence and a restricted proof of concept. & Evidence and responsible owner to be recorded \\
5 & Define acceptance criteria for identity, semantics, provenance, latency, resilience, and governance. & Evidence and responsible owner to be recorded \\
6 & Record which standards are implementation tools and which local mappings and controls remain necessary. & Evidence and responsible owner to be recorded \\
\end{longtable}

\subsection{Illustrative Outputs}
The outputs in Table~\ref{tab:comp-ch02-outputs} show how technical,
clinical, operational, economic, equity, and governance findings should be kept
visible rather than compressed into a single favourable or unfavourable score.

\begin{longtable}{@{}p{0.30\textwidth}p{0.24\textwidth}p{0.38\textwidth}@{}}
\caption{Illustrative model and decision outputs}
\label{tab:comp-ch02-outputs}\\
\toprule
\textbf{Output} & \textbf{Result} & \textbf{Interpretation} \\
\midrule
\endfirsthead
\toprule
\textbf{Output} & \textbf{Result} & \textbf{Interpretation} \\
\midrule
\endhead
\bottomrule
\endlastfoot
Immediate procurement feasibility & Low & Critical data and identity dependencies unresolved \\
Foundation reuse potential & High & Seven planned services share the same capabilities \\
Short-term delay cost & Moderate & Potential benefit deferred for nine months \\
Risk of isolated procurement & High & Repeated mapping, weak provenance, and technical debt \\
Preferred decision & Foundation-first with restricted proof of concept & Allows learning without operational dependence \\
\end{longtable}

\subsection{Sensitivity, Uncertainty, and Subgroups}
At minimum, the completed analysis should vary the parameters most likely to
change the decision, test at least one credible alternative model structure, and
report subgroup results separately where performance, access, response, burden,
or opportunity cost may differ. Deterministic analysis should identify thresholds;
probabilistic analysis should be used where credible parameter distributions and
correlations are available. Scenario analysis is preferable when structural or
future uncertainty cannot be represented defensibly by one probability model.

The record should distinguish uncertainty that can be reduced through evidence,
uncertainty that can be managed through safeguards, and uncertainty that remains
too consequential to accept. Missing evidence should not be represented as an
average or neutral value without justification.

\subsection{Interpretation and Recommendation}
Fund the shared digital-foundation programme and permit a restricted, non-clinically dependent proof of concept in one service. Procurement for routine use should wait until predefined data, interoperability, resilience, and governance thresholds are met.

The recommendation is conditional rather than permanent. It applies only to the
specified population, setting, intervention configuration, evidence cut-off, and
version. The following conditions should be incorporated into the approval,
contract, implementation, monitoring, or renewal record:

\begin{itemize}
 \item identity and duplicate-record controls meet the agreed threshold.
 \item medication and discharge data have documented provenance and completeness.
 \item interfaces meet decision-latency and recovery requirements.
 \item local FHIR profiles or equivalent exchange specifications and analytical transformations are version controlled.
 \item the investment case identifies portfolio reuse and avoids assigning all shared cost to one application.
\end{itemize}

\subsection{Decision Statement}
\begin{quote}
For the specified decision problem and comparator set, the institution should
\emph{[proceed / proceed with conditions / redesign / pilot / restrict / defer /
replace / retire]}. The decision applies to \emph{[population, setting, version,
and evidence period]}, is owned by \emph{[accountable authority]}, and will be
reviewed on \emph{[date]} or earlier if \emph{[trigger events]} occur. The
evidence, assumptions, unresolved uncertainty, mandatory safeguards, monitoring
thresholds, and exit arrangements are recorded in the accompanying dossier.
\end{quote}
